\section{Possible systems, evidence and action}\label{sec:foundations}
\paragraph{The set from which a claim follows.}
A \emph{completion} \(C\) specifies the relevant system semantics: which states, observations, actions, transitions and effects are possible. A probabilistic completion also specifies their conditional laws. A physical detail can be omitted if the specification and the available interfaces do not depend on it. Fix the evidence \(\Omega\), including its independently warranted interpretation. Write \(C\models\Omega\) when \(C\) satisfies it. Fix a target map \(\Psi\) before examining the audit outcome. Its value might be a Boolean guarantee, a cost or a vector of possible intervention costs. Define
\begin{equation}
 {\cal C}_{\Omega}=\{C:C\models\Omega\},\qquad
 {\cal I}_{\Psi}(\Omega)=\{\Psi(C):C\in{\cal C}_{\Omega}\}.
 \label{eq:identified}
\end{equation}
The first set lists the systems we still take seriously. The second lists their answers. In Figure~\ref{fig:worlds}, after the good guard report the first set has two members and the second has one:
\[
 |{\cal C}_\Omega|>1,\qquad |{\cal I}_\Psi(\Omega)|=1.
\]
Here the target is identified while the system remains unknown. An empty completion class signals inconsistent evidence or assumptions. A singleton \emph{answer set} identifies the target. Two or more values leave the target ambiguous. An identified answer can be negative. Sound use in deployment separately requires the actual installation to belong to the class. This membership is a source claim with its own evidence.

\paragraph{How a read changes the set.}
An observation is a record delivered to the controller. A trusted source fact has an established link to the particular installation, field and meaning it reports. A pure, exact read \(j\) with report function \(r_j\) and observed value \(z\) intersects the class with \(\{C:r_j(C)=z\}\). Permissions, possible reports and prices belong to the \emph{query interface}. Installing a new guard can change \(C\) itself. It requires an intervention model relating the original and resulting installations. A receipt for the change cannot establish that the original installation had the new property.

\paragraph{What an implementable controller knows.}
A history contains the controller's previous actions, delivered observations and available receipts. A policy assigns actions or distributions over them to histories. Nonanticipativity means that equal observable histories induce the same decision law. Hidden state, future draws and the evaluator's separate ground-truth record are unavailable to the policy. Every action it can select must have authority in each still-possible physical state in which selection can occur. For a finite sure guarantee we use a proper policy tree, one whose supported branches all terminate. We count the entire positive support, meaning every outcome with strictly positive probability or every declared possible outcome of a nondeterministic command. An outcome with a small positive probability still counts for a sure guarantee.

\paragraph{The original closure problem.}
The original Perceptive Closure manuscript and theorem charter ask whether a repair would still work if one of its information resources were held fixed \citep{d35,charter}.
At a checkpoint a system may have accumulated evidence, chosen how to interpret that evidence, and obtained certain permissions. Perceptive Closure treats these as distinct resources. Evidence \(E\) is the admitted history representative. Representation \(R\) tells the certificate which history representatives count as equivalent. Authority \(A\) tells it which capabilities are admissible. Ordinary preconditions, such as the current workflow phase, also constrain execution. The controller's observation can be finer or coarser than the certificate's representation.

In one fixed contract let \(s\) denote a checkpoint, \(H(s)\) its admitted representative, \(P(s)\) the equivalence relation there, and \(\Lambda(s)\) its effective authority. Let \(x\) denote a latent world in that contract, \(u\) a possible history representative and \(x\sim u\) their compatibility relation. The worlds represented at the checkpoint are
\[
 C(s)=\{x:\ \text{some }u\in[H(s)]_{P(s)}\text{ satisfies }x\sim u\}.
\]
We assume these sets are nonempty on reachable checkpoints. Fix the local target label \(a(x)\). The checkpoint is \emph{closed} when all represented worlds have the same label. Uncertainty over \(x\) inside one contract and uncertainty over the contract \(C\) in \eqref{eq:identified} are separate. A policy that stops only at physically closed checkpoints can still need an extra read if its task requires it to \emph{report} a common label across those checkpoints.

A \emph{source-atomic} operation is a whole command supplied by the actual interface. To decide whether it touches \(E,R,A\), take the union of the changes in \(H,P,\Lambda\) over all of its supported successors, with equality of representations and permissions judged by their effective behavior. Freezing a resource deletes the entire operation at any checkpoint where it could change that resource. A caller cannot keep a favorable branch of an indivisible command. For a set \(F\) of frozen resources, let \(\kappa_F\) be the minimum worst supported cost of a proper closing policy, and \(+\infty\) if none exists. Its \emph{signature} consists of the eight \(\kappa_F\) in the order
\[
 \varnothing,\ E,\ R,\ A,\ ER,\ EA,\ RA,\ ERA,
\]
where juxtaposition denotes a set of resources. A successful unrestricted run records one path cost. Establishing even the first coordinate requires a policy bound on all supported branches and a matching lower bound. Appendix~\ref{app:foundations} derives the consequences of this distinction.
